\chapter{Eksponen dan Logaritma}\label{ch:g12:exp}

\omterm{thm:g12:exp:existence}{Fungsi eksponen} adalah satu-satunya \omterm{def:g11:func:function}{fungsi} yang sama dengan \omterm{def:g12:deriv:derivative}{turunannya} sendiri
dan bernilai $1$ di $0$. \omterm{def:g11:func:function}{Fungsi} itu mengubah jumlah menjadi hasil kali;
inversnya, \omterm{def:g12:exp:ln}{logaritma asli}, mengubah hasil kali menjadi jumlah. Bersama-sama
keduanya memerikan setiap gejala yang laju perubahannya sebanding dengan
ukurannya: peluruhan radioaktif, pertumbuhan populasi, bunga majemuk.

\section{Fungsi eksponen}

\begin{theorem}[Keberadaan dan ketunggalan]\label{thm:g12:exp:existence}
Ada satu \omterm{def:g11:func:function}{fungsi} $f \colon \R \to \R$ yang tunggal dan \omterm{def:g12:deriv:derivative}{dapat diturunkan} sehingga
\[
f' = f \qquad\text{dan}\qquad f(0) = 1 .
\]
\omterm{def:g11:func:function}{Fungsi} itu disebut \emph{fungsi eksponen}\index{eksponen} dan ditulis
$\exp$, atau $x \mapsto \eu^x$.
\end{theorem}

\begin{proof}[Bukti ketunggalannya]
Pertama, \omterm{def:g11:func:function}{fungsi} semacam itu tidak pernah lenyap. Memang, misalkan
$g(x) = f(x)f(-x)$; maka
\[
g'(x) = f'(x)f(-x) - f(x)f'(-x) = f(x)f(-x) - f(x)f(-x) = 0,
\]
sehingga $g$ tetap dan sama dengan $g(0) = 1$: jadi untuk setiap $x$ berlaku
$f(x)f(-x) = 1$, dan khususnya $f(x) \neq 0$.

Kini misalkan $f_1, f_2$ dua penyelesaian dan $h = \dfrac{f_1}{f_2}$ (sah
sebab $f_2$ tidak pernah lenyap). Maka
\[
h' = \frac{f_1' f_2 - f_1 f_2'}{f_2^2} = \frac{f_1 f_2 - f_1 f_2}{f_2^2} = 0,
\]
sehingga $h$ tetap dan sama dengan $h(0) = 1$, \emph{yaitu} $f_1 = f_2$.

\emph{Keberadaannya} diterima tanpa bukti pada tingkat ini (dapat diperoleh
lewat \omterm{def:g12:seq:sequence}{barisan} pada \cref{exo:g12:seq:10}, atau sebagai invers \omterm{def:g12:exp:ln}{logaritma}
yang disusun lewat pengintegralan pada \cref{ch:g12:integ}).
\end{proof}

\begin{proposition}[Persamaan fungsional]\label{prop:g12:exp:funceq}
Untuk semua $x, y \in \R$ dan $n \in \Z$:
\[
\eu^{x+y} = \eu^x\,\eu^y, \qquad
\eu^{-x} = \frac{1}{\eu^x}, \qquad
\eu^{x-y} = \frac{\eu^x}{\eu^y}, \qquad
\bigl(\eu^{x}\bigr)^n = \eu^{nx}.
\]
Lebih lanjut, $\eu^x > 0$ untuk semua $x$.
\end{proposition}

\begin{proof}
Tetapkan $y$ lalu tinjaulah $\varphi(x) = \dfrac{\exp(x+y)}{\exp(x)\exp(y)}$.
Pembilang dan penyebutnya, sebagai \omterm{def:g11:func:function}{fungsi} dari $x$, sama-sama merupakan
penyelesaian $f' = f$ sampai pada konstanta $\exp(y)$; menurunkan $\varphi$
secara langsung (aturan hasil bagi) memberi $\varphi' = 0$, sehingga
$\varphi \equiv \varphi(0) = 1$, dan itu membuktikan identitas pertamanya.
Mengambil $y = -x$ memberi yang kedua (dengan nilai $\eu^0 = 1$), lalu yang
ketiga menyusul. Yang keempat adalah induksi dari yang pertama untuk
$n \geq 0$, yang diperluas ke $n < 0$ lewat yang kedua.

Kepositifannya: $\eu^x = \left(\eu^{x/2}\right)^2 \geq 0$ dan $\eu^x \neq 0$
(ditunjukkan pada \cref{thm:g12:exp:existence}), sehingga $\eu^x > 0$.
\end{proof}

\begin{proposition}[Variasi dan limit]\label{prop:g12:exp:variations}
\omterm{thm:g12:exp:existence}{Fungsi eksponen} tegas naik, \omterm{def:g12:deriv:convex}{cembung}, dan
\[
\lim_{x\to-\infty} \eu^x = 0, \qquad
\lim_{x\to+\infty} \eu^x = +\infty .
\]
\end{proposition}

\begin{proof}
$(\exp)' = \exp > 0$ memberi kenaikan yang tegas; $(\exp)'' = \exp > 0$
memberi \omterm{def:g12:deriv:convex}{kecembungannya}. Menurut \omterm{def:g12:deriv:convex}{kecembungannya}, $\eu^x \geq 1 + x$ (\omterm{def:g12:deriv:tangent}{garis
singgung} di $0$, lihat \cref{ex:g12:deriv:ineq}), sehingga $\eu^x \to +\infty$
ketika $x \to +\infty$ lewat pembandingan. Lalu $\eu^{x} = 1/\eu^{-x} \to 0$
ketika $x \to -\infty$.
\end{proof}

\begin{omfigure}
\begin{tikzpicture}
\begin{axis}[omaxis,
  width=8.8cm, height=6cm,
  xmin=-3.3, xmax=2.4, ymin=-2.4, ymax=7.8,
  xtick={-2,-1,1,2}, ytick={1}]
  \addplot[omProp, thick, domain=-3.2:2.3] {1 + x};
  \addplot[omDef, thick, domain=-3.2:2.0] {exp(x)};
  \fill (axis cs:0,1) circle (1.5pt);
  \node[omDef, font=\small, anchor=east] at (axis cs:1.35,5.2) {$\eu^x$};
  \node[omProp, font=\small, anchor=north west] at (axis cs:1.05,1.6) {$y = 1+x$};
\end{axis}
\end{tikzpicture}

{\small \omterm{thm:g12:exp:existence}{Fungsi eksponen}: tegas naik, \omterm{def:g12:deriv:convex}{cembung}, dengan
$\lim_{-\infty} \eu^x = 0$. \omterm{def:g11:func:function}{Fungsi} itu terletak di atas \omterm{def:g12:deriv:tangent}{garis singgungnya} di
$0$: $\eu^x \geq 1 + x$.}
\end{omfigure}

\begin{theorem}[Pembandingan pertumbuhan]\label{thm:g12:exp:growth}
\index{pembandingan pertumbuhan}
Untuk setiap \omterm{def:g10:numbers:sets}{bilangan bulat} $n \geq 1$,
\[
\lim_{x\to+\infty} \frac{\eu^x}{x^n} = +\infty,
\qquad
\lim_{x\to-\infty} x^n \eu^x = 0 .
\]
Dengan kata lain: \omterm{thm:g12:exp:existence}{fungsi eksponen} mengalahkan setiap pangkat $x$.
\end{theorem}

\begin{proof}
Untuk $n = 1$: menerapkan $\eu^t \geq 1 + t \geq t$ pada $t = x/2$,
\[
\frac{\eu^x}{x} = \frac{\left(\eu^{x/2}\right)^2}{x}
\geq \frac{(x/2)^2}{x} = \frac{x}{4} \xrightarrow[x\to+\infty]{} +\infty .
\]
Untuk $n$ yang umum, tulislah
\[
\frac{\eu^x}{x^n}
= \frac{1}{n^n}\left(\frac{\eu^{x/n}}{x/n}\right)^{n}
\xrightarrow[x\to+\infty]{} +\infty
\]
menurut kasus $n=1$ dan komposisinya. Limitnya di $-\infty$ menyusul lewat
substitusi $x \mapsto -x$:
$\abs{x^n \eu^x} = \frac{\abs{x}^n}{\eu^{\abs x}} \to 0$.
\end{proof}

\begin{omfigure}
\begin{tikzpicture}
\begin{axis}[omaxis,
  width=8.8cm, height=6cm,
  xmin=-0.4, xmax=5.6, ymin=-12, ymax=155,
  xtick={1,2,3,4,5}, ytick={50,100,150}]
  \addplot[omThm, thick, domain=0:5.4] {x^2};
  \addplot[omProp, thick, domain=0:5.3] {x^3};
  \addplot[omDef, thick, domain=-0.4:5.0] {exp(x)};
  \fill (axis cs:4.536,93.4) circle (1.5pt);
  \draw[black, dashed, thin] (axis cs:4.536,0) -- (axis cs:4.536,93.4);
\end{axis}
\end{tikzpicture}

{\small \omterm{thm:g12:exp:existence}{Fungsi eksponen} (biru) akhirnya mengalahkan setiap pangkat: \omterm{def:g11:func:function}{fungsi} itu
menyalip $x^2$ (merah) sejak awal dan $x^3$ (jingga) di $x \approx 4.54$.}
\end{omfigure}

\section{Logaritma asli}

\begin{definition}[Logaritma asli]\label{def:g12:exp:ln}
\omterm{thm:g12:exp:existence}{Fungsi eksponen} \omterm{def:g12:limcont:continuity}{kontinu} dan tegas naik dari $\R$ pada
$\intoo{0}{+\infty}$; menurut teorema bijeksi
(\cref{thm:g12:limcont:bijection}), untuk setiap $y > 0$ \omterm{def:g10:algebra:equation}{persamaan}
$\eu^x = y$ mempunyai penyelesaian tunggal. Penyelesaian itu disebut
\emph{logaritma asli}\index{logaritma} dari $y$, ditulis $\ln y$. Jadi
\[
\text{untuk } x \in \R,\ y > 0: \qquad y = \eu^x \iff x = \ln y .
\]
Khususnya $\ln 1 = 0$, $\ln \eu = 1$, serta
$\eu^{\ln y} = y$ dan $\ln(\eu^x) = x$.
\end{definition}

\begin{omfigure}
\begin{tikzpicture}
\begin{axis}[omaxis,
  width=8.6cm, axis equal image,
  xmin=-4.3, xmax=6.3, ymin=-4.3, ymax=6.3,
  xtick={1}, ytick={1}]
  \addplot[black, dashed, domain=-3.9:5.9] {x};
  \addplot[omDef, thick, domain=-4.1:1.79] {exp(x)};
  \addplot[omThm, thick, domain=0.0167:5.9, samples=120] {ln(x)};
  \node[omDef, font=\small, anchor=east] at (axis cs:0.95,2.6) {$\eu^x$};
  \node[omThm, font=\small, anchor=north] at (axis cs:4.3,1.25) {$\ln x$};
  \node[font=\small, anchor=west, rotate=45] at (axis cs:4.6,4.85) {$y=x$};
\end{axis}
\end{tikzpicture}

{\small Kurva $\exp$ dan $\ln$ saling mencerminkan
terhadap garis $y = x$: titik $(x, \eu^x)$ tercermin menjadi $(\eu^x, x)$.}
\end{omfigure}

\begin{proposition}[Sifat aljabar]\label{prop:g12:exp:lnalg}
Untuk semua $a, b > 0$ dan $n \in \Z$:
\[
\ln(ab) = \ln a + \ln b, \quad
\ln\frac{1}{a} = -\ln a, \quad
\ln\frac{a}{b} = \ln a - \ln b, \quad
\ln(a^n) = n \ln a, \quad
\ln\sqrt{a} = \tfrac12 \ln a .
\]
\end{proposition}

\begin{proof}
$\eu^{\ln a + \ln b} = \eu^{\ln a}\eu^{\ln b} = ab$, lalu mengambil $\ln$
kedua ruasnya memberi identitas pertamanya. Yang lain menyusul lewat mekanisme
yang sama dari identitas \cref{prop:g12:exp:funceq} yang bersesuaian.
\end{proof}

\begin{proposition}[Sifat analitis]\label{prop:g12:exp:lnanalytic}
\omterm{def:g11:func:function}{Fungsi} $\ln$ \omterm{def:g12:deriv:derivative}{dapat diturunkan} pada $\intoo{0}{+\infty}$ dengan
\[
(\ln x)' = \frac{1}{x},
\]
tegas naik, \omterm{def:g12:deriv:convex}{cekung}, serta
$\lim\limits_{x\to0^+} \ln x = -\infty$ dan
$\lim\limits_{x\to+\infty} \ln x = +\infty$. Lebih lanjut, untuk setiap
\omterm{def:g10:numbers:sets}{bilangan bulat} $n \geq 1$,
\[
\lim_{x\to+\infty} \frac{\ln x}{x^{1/n}} = 0, \qquad
\lim_{x\to0^+} x \ln x = 0 .
\]
\end{proposition}

\begin{proof}
Keterdiferensialan \omterm{def:g11:func:function}{fungsi} inversnya diterima tanpa bukti pada tingkat ini;
dengan menerimanya, turunkan identitas $\eu^{\ln x} = x$ lewat aturan rantai:
$(\ln x)'\,\eu^{\ln x} = 1$, sehingga $(\ln x)' = \frac{1}{x} > 0$, dan dari
sana kenaikan yang tegas, lalu $(\ln x)'' = -\frac1{x^2} < 0$, dan dari sana
kecekungannya. Limitnya di $0^+$ dan $+\infty$ mencerminkan limit $\exp$
lewat bijeksinya. Untuk \omterm{thm:g12:exp:growth}{pembandingan pertumbuhannya}, substitusikan
$x = \eu^{t}$: $\frac{\ln x}{x} = \frac{t}{\eu^t} \to 0$ ketika
$t \to +\infty$ menurut \cref{thm:g12:exp:growth}, dan serupa untuk limit
lainnya dengan $x = \eu^{-t}$, yang memberi $x\ln x = -t\eu^{-t} \to 0$.
\end{proof}

\begin{method}[Menyelesaikan persamaan dengan $\exp$ dan $\ln$]\label{met:g12:exp:equations}
Kedua \omterm{def:g11:func:function}{fungsinya} tegas naik, sehingga keduanya dapat diterapkan pada (atau
dihapus dari) kedua ruas sebuah \omterm{def:g10:algebra:equation}{persamaan} atau pertidaksamaan tanpa mengubah arahnya:
\[
\eu^{u} = \eu^{v} \iff u = v, \qquad
\eu^{u} \leq \eu^{v} \iff u \leq v,
\]
dan demikian pula $\ln$ pada besaran yang positif. \emph{Periksalah selalu
\omterm{def:g11:func:function}{daerah asalnya} lebih dahulu} ($\ln$ menuntut argumen yang positif). Untuk
\omterm{def:g10:algebra:equation}{persamaan} seperti $a^x = b$ dengan $a>0$, $a \neq 1$, tulis ulang
$a^x = \eu^{x\ln a}$ lalu selesaikan $x = \frac{\ln b}{\ln a}$.
\end{method}

\begin{example}\label{ex:g12:exp:halflife}
Sebuah cuplikan radioaktif meluruh mengikuti $N(t) = N_0\,\eu^{-\lambda t}$.
\emph{Waktu paruhnya} $T$ memenuhi $N(T) = N_0/2$, \emph{yaitu}
$\eu^{-\lambda T} = \frac12$, sehingga $T = \frac{\ln 2}{\lambda}$: jadi waktu
paruhnya tidak bergantung pada banyaknya bahan mula-mula.
\end{example}

\section{Latihan}

\begin{exercise}[$\star$]\label{exo:g12:exp:1}
Sederhanakan
$\dfrac{\eu^{3x}\,\eu^{-x+1}}{\eu^{x}}$ dan
$\ln\!\left(\dfrac{\eu^2\sqrt{\eu}}{\eu^{-3}}\right)$.
\end{exercise}

\begin{exercise}[$\star$]\label{exo:g12:exp:2}
Selesaikan dalam $\R$:
\[
\text{(a) } \eu^{2x} - 3\eu^{x} + 2 = 0;
\qquad
\text{(b) } \ln(x - 1) + \ln(x + 2) = \ln 4 .
\]
\end{exercise}

\begin{exercise}[$\star$]\label{exo:g12:exp:3}
Hitunglah limitnya:
\[
\lim_{x\to+\infty} \frac{\eu^x - x^2}{\eu^x + 1}, \qquad
\lim_{x\to+\infty} \frac{\ln x}{\sqrt x}, \qquad
\lim_{x\to0^+} x^2 \ln x, \qquad
\lim_{x\to+\infty} \bigl(x - \ln x\bigr).
\]
\end{exercise}

\begin{exercise}[$\star\star$]\label{exo:g12:exp:4}
Telaahlah \omterm{def:g11:func:function}{fungsi} $f(x) = x\,\eu^{-x}$ pada $\R$: variasi, limit, dan \omterm{def:g12:deriv:extremum}{nilai
ekstremnya}; lalu tunjukkan bahwa kurvanya mempunyai sebuah \omterm{def:g12:deriv:inflection}{titik belok} dan
berikan \omterm{def:g10:coordgeom:system}{koordinatnya}.
\end{exercise}

\begin{exercise}[$\star\star$]\label{exo:g12:exp:5}
Tunjukkan bahwa untuk semua $x > -1$ berlaku $\ln(1 + x) \leq x$, dengan
kesamaan hanya di $x = 0$. Simpulkan bahwa untuk semua $n \geq 1$,
\[
\left(1 + \frac{1}{n}\right)^n \leq \eu \leq \left(1 - \frac{1}{n+1}\right)^{-(n+1)} .
\]
\end{exercise}

\begin{exercise}[$\star\star$]\label{exo:g12:exp:6}
Sebuah modal $C_0$ ditanamkan dengan laju $3\%$ per tahun, dengan bunga
dimajemukkan setiap tahun.
\begin{enumerate}
  \item Nyatakan modalnya $C_n$ setelah $n$ tahun.
  \item Setelah berapa tahun modalnya berlipat dua? Berikan jawaban eksaknya
        dengan $\ln$, lalu sebuah nilai numeris.
  \item Bandingkan dengan \omterm{def:g10:numbers:approx}{hampiran} ``$70$ dibagi lajunya dalam
        persen'' yang dipakai para bankir.
\end{enumerate}
\end{exercise}

\begin{exercise}[$\star\star$]\label{exo:g12:exp:7}
Selesaikan pertidaksamaan $\eu^{2x} - \eu^{x+1} > 0$, lalu pertidaksamaan
$\ln(x^2 - 1) \leq \ln(x + 5)$.
\end{exercise}

\begin{exercise}[$\star\star\star$]\label{exo:g12:exp:8}
Misalkan $f(x) = \dfrac{\eu^x}{x}$ untuk $x > 0$.
\begin{enumerate}
  \item Telaahlah variasi $f$ pada $\intoo{0}{+\infty}$ lalu berikan
        \omterm{def:g10:functions:extrema}{minimumnya}.
  \item Untuk nilai $k$ yang mana \omterm{def:g10:algebra:equation}{persamaan} $\eu^x = kx$ mempunyai $0$, $1$
        atau $2$ penyelesaian di $\intoo{0}{+\infty}$?
\end{enumerate}
\end{exercise}

\begin{exercise}[$\star\star\star$]\label{exo:g12:exp:9}
Untuk $n \geq 1$, misalkan $u_n = \left(1 + \frac1n\right)^n$.
\begin{enumerate}
  \item Dengan \cref{exo:g12:exp:5}, tunjukkan bahwa $(u_n)$ \omterm{def:g12:seq:bounded}{terbatas di atas}
        oleh~$\eu$.
  \item Tunjukkan bahwa $\ln u_n = n \ln\left(1 + \frac1n\right) \to 1$, lalu
        simpulkan bahwa $u_n \to \eu$.
        (Petunjuk: kenalilah sebuah hasil bagi selisih $\ln$ di $1$.)
\end{enumerate}
\end{exercise}

\section{Soal: Logaritma menjinakkan dunia}

\begin{problem}[{Soal akhir pekan --- waktu pelipatduaan dan kaidah
72, skala gempa bumi, asam dan piano, serta tetapan $\eu$ yang
meninggalkan sidik jari di mana-mana}]
\label{pb:g12:exp:1}
Apa pun yang tumbuh dengan \emph{persentase} yang tetap tumbuh secara
eksponensial --- tabungan, bakteri, wabah --- dan apa pun yang membentang
melewati terlalu banyak pangkat sepuluh untuk digenggam --- energi gempa bumi,
keasaman, intensitas bunyi --- dijinakkan oleh \omterm{def:g12:exp:ln}{logaritma}. Soal ini
menghitung waktu pelipatduaan lalu menyingkap kaidah 72 milik para bankir,
membaca skala logaritmis dunia, lalu mengumpulkan sidik jari yang ditinggalkan
bilangan $\eu$ di setiap tempat \omterm{def:g11:prob:model}{kejadian} perkara
(\cref{prop:g12:exp:funceq}, \cref{thm:g12:exp:growth},
\cref{met:g12:exp:equations}).

\medskip
\noindent\textbf{Bagian I --- Kefasihan.}
\begin{enumerate}
  \item Selesaikan: $\eu^{2x} = 5$; \quad $\ln(3x - 1) = 2$; \quad
        $\eu^{2x} - 3\eu^x + 2 = 0$ (bentuk kuadrat dalam samaran).
  \item Sederhanakan: $\dfrac{\ln 8}{\ln 2}$; \quad
        $\ln\!\left(\eu^3 \sqrt{\eu}\right)$; \quad
        $\eu^{\ln 5 - \ln 2}$.
  \item Buktikan ketaksamaan induknya: $\eu^x \geq 1 + x$ untuk setiap
        \omterm{def:g10:numbers:sets}{bilangan real} $x$ (telaahlah $f(x) = \eu^x - 1 - x$), lalu simpulkan
        $\ln(1 + u) \leq u$ untuk $u > -1$.
  \item Hitunglah
        $\lim_{x \to +\infty} x^2 \eu^{-x}$,
        $\lim_{x \to +\infty} \frac{\ln x}{x}$
        (\cref{thm:g12:exp:growth}), dan
        $\lim_{x \to 0} \frac{\eu^x - 1}{x}$ (kenalilah sebuah
        \omterm{def:g12:deriv:derivative}{turunan}).
  \item Telaahlah $f(x) = x \ln x$ pada $\intoo{0}{+\infty}$:
        variasinya, \omterm{def:g10:functions:extrema}{minimumnya}, dan limitnya di $0^+$ (diterima tanpa bukti:
        $x \ln x \to 0$). \omterm{def:g11:func:function}{Fungsi} kecil ini mengukur
        informasi dan entropi di sepanjang jilid universitas.
\end{enumerate}

\medskip
\noindent\textbf{Bagian II --- Waktu pelipatduaan dan kaidah
72.}
\begin{enumerate}[resume]
  \item Tabungan tumbuh $3\,\%$ per tahun. Selesaikan
        $1.03^n = 2$: berapa lama modalnya berlipat dua?
  \item Para bankir menaksir waktu pelipatduaan sebagai
        $\frac{72}{\text{laju dalam }\%}$. Ujilah kaidah itu pada
        $3\,\%$, $6\,\%$ dan $9\,\%$ terhadap nilai eksak
        $\frac{\ln 2}{\ln(1 + r)}$, lalu jelaskan: untuk $r$ yang
        kecil, $\ln(1 + r) \approx r$ (ketaksamaan pertanyaan 3 adalah
        separuh ceritanya), sehingga tetapan eksaknya adalah
        $100 \ln 2 \approx 69.3$ --- lalu mengapa para bankir lebih menyukai
        $72$?
  \item Sebuah bakteri membelah setiap $20$ menit:
        $p(t) = 2^{t/20}$ ($t$ dalam menit). Tulis ulang sebagai
        $\eu^{\lambda t}$, lalu hitunglah $p$ setelah $24$ jam.
        Jawabannya (lebih dari $10^{21}$) membuktikan apa tentang
        modelnya --- dan apa yang menghentikan koloni yang sesungguhnya?
  \item Kafeina meninggalkan tubuh dengan waktu paruh sekitar $5$
        jam. Setelah secangkir kopi berkadar $100$ mg pada pukul 15.00, berapa
        yang tersisa pada pukul 23.00? Pada pukul berapa kadarnya turun di
        bawah $10$ mg? (Susah tidur pun berlogaritma.)
  \item Satu euro dengan bunga $100\,\%$ per tahun: hitunglah
        modal akhir tahunnya bila bunganya dimajemukkan setiap tahun, setiap
        bulan dan setiap hari, lalu berikan langit-langit yang dibuktikan tak
        tertembus oleh \cref{exo:g12:exp:9}. Tetapan termasyhur mana
        yang menjadi limit ketamakan murni?
  \item Mengapa para ilmuwan menggambar $\ln(\text{cacah kasus})$ terhadap
        waktu pada tahap awal sebuah wabah? Apa yang disingkapkan sebuah
        garis lurus pada gambar itu, dan apa yang diukur \omterm{def:g10:reffunc:affine}{gradiennya}?
\end{enumerate}

\medskip
\noindent\textbf{Bagian III --- Skala logaritmis dunia.}
(Tulislah $\log_{10} x = \frac{\ln x}{\ln 10}$.)
\begin{enumerate}[resume]
  \item Taraf bunyi dalam desibel: $L = 10 \log_{10}(I/I_0)$. Sebuah
        percakapan terukur $60$ dB, sebuah konser rok $120$ dB:
        berapa kali lipat intensitas bunyinya berbeda?
  \item Magnitudo gempa bumi naik $1$ ketika
        amplitudo seismiknya dikalikan $10$, dan energi
        yang dilepaskannya berskala seperti amplitudo$^{3/2}$. Bandingkan
        gempa bermagnitudo $5$ dan bermagnitudo $7$: perbandingan
        amplitudonya, lalu perbandingan energinya.
  \item Kimia: $\text{pH} = -\log_{10}[\mathrm{H^+}]$.
        Air jeruk lemon ber-pH $2$, susu ber-pH $6.5$: berapa
        perbandingan kepekatan asamnya?
  \item Musik: setiap oktaf melipatduakan frekuensinya. Dari
        A terendah pada piano ($27.5$ Hz) sampai C tertingginya
        ($4\,186$ Hz), berapa oktaf yang dibentangkan papan tutsnya?
        Lalu mengapa indra kita --- pendengaran, penglihatan, perasaan
        terhadap gempa --- lebih menyukai skala logaritmis? (Satu kalimat.)
  \item Mistar hitung, kalkulator para insinyur selama $350$
        tahun: dua batang bertanda sehingga \emph{panjang}
        sampai tanda $x$ sebanding dengan $\ln x$. Jelaskan bagaimana
        menggeser satu batang di sepanjang batang yang lain mengalikan
        bilangan, lalu sebutkan identitas
        \cref{prop:g12:exp:lnalg} yang melakukan pekerjaannya.
\end{enumerate}

\medskip
\noindent\textbf{Bagian IV --- Sidik jari $\eu$.}
\begin{enumerate}[resume]
  \item Dari \cref{exo:g12:exp:8}: \omterm{def:g10:algebra:equation}{persamaan}
        $\eu^x = kx$ ($x > 0$) mempunyai $0$, $1$ atau $2$ penyelesaian
        menurut kedudukan $k$ terhadap sebuah
        ambang. Nyatakan ulang hasilnya --- lalu periksalah fakta
        geometris di baliknya: garis $y = \eu x$ tepat merupakan \omterm{def:g12:deriv:tangent}{garis
        singgung} \omterm{thm:g12:exp:existence}{fungsi eksponen} yang melalui titik asal.
  \item Tetapan nyaris-meleset: dengan $n$ karcis undian yang
        masing-masing berpeluang menang $\frac1n$,
        $\P(\text{tidak menang}) = \left(1 - \frac1n\right)^n$.
        Hitunglah limitnya lewat
        $n \ln\!\left(1 - \frac1n\right)$ (pakailah limit baku
        $\frac{\ln(1 + u)}{u} \to 1$), lalu cocokkan dengan
        $0.368$ yang misterius pada \cref{pb:g11:binom:1}.
  \item Kaidah $37\,\%$: untuk memilih yang terbaik di antara $n$
        calon yang diwawancarai dalam urutan acak (tanpa boleh kembali),
        strategi optimumnya menolak
        $\frac{n}{\eu} \approx 37\,\%$ yang pertama lalu mengambil calon
        pertama yang lebih baik daripada semua calon sebelumnya --- dan
        berhasil dengan peluang sekitar $\frac1\eu$. Dari mana kedua
        $\frac1\eu$ pada pertanyaan 18 dan 19 berasal ---
        nyatakan mekanisme bersamanya (banyak \omterm{def:g11:prob:model}{kejadian} saling bebas yang
        berpeluang kecil), lalu hitunglah $\frac1\eu$ sampai tiga
        angka desimal.
  \item Penutup --- potret $\eu$: langit-langit pemajemukan
        (pertanyaan 10); bilangan pokok yang \omterm{def:g12:deriv:tangent}{garis singgungnya} di $0$
        bergradien tepat $1$ (pertanyaan 3 dan 4); pertumbuhan yang tidak
        terkejar polinomial mana pun (pertanyaan 4); tetapan
        nyaris-meleset dan penghentian optimum (pertanyaan 18--19);
        serta inversnya $\ln$, yang mengubah hasil kali menjadi jumlah
        (pertanyaan 16) dan dasawarsa menjadi inci (Bagian III). Masing-masing
        satu kalimat.

\end{enumerate}
\end{problem}
